\section{Introduction}
Message passing describes computation through entities, interactions, local
functions and reductions. In machine learning, this pattern includes graph
neural networks, point-cloud neighborhoods and attention; numerical simulation
uses related structures for particles and mesh operators. The same high-level
pattern can require very different execution strategies. A stored sparse graph
needs indexed traversal, a radius relation can exploit spatial locality, and
a dense relation need not have an explicit edge array at all.

Tensor-oriented implementations often build connectivity first, gather fields
for every edge, compute messages, and finally reduce them. For edge count
and feature width denoted by $E$ and $F$, a materialized message array uses
storage proportional to $EF$, even when its sole consumer produces $NF$
output values for $N$ destinations. Dynamic coordinates can make the cost of
constructing these intermediates recur at every call. When the input graph
itself exceeds device memory, fusion alone is insufficient: execution must
also control which topology and field pages reside on the accelerator.

Tiga, short for \emph{Target-Independent Graph Acceleration}, separates the
meaning of a message-passing program from its traversal, hardware mapping and
storage residence. Target independence describes this separation, not equal
backend coverage: current executable routes target CPUs and NVIDIA GPUs.
The central design choice is to retain relation and reducer information before
lowering to tensor and device operations. This enables two complementary
mechanisms: consuming generated neighbors directly, and streaming stored
graphs with a bounded temporary working set.

This paper contributes a relation-aware compilation interface and lowering
strategy, a destination-page execution contract for graphs exceeding device
memory, and an evaluation of their forward latency, allocation and capacity.
The results establish useful execution paths rather than a universal framework
ranking. In particular, the evaluation does not isolate every compiler
optimization or establish end-to-end training throughput. Differentiation and
distributed execution are supported on specific paths, with their boundaries
stated separately from the central performance claims.

\paragraph{Relation to prior systems.}
PyTorch Geometric exposes message passing to graph-learning applications
\citep{pyg,pyg2}. Graphiler retains graph-related dataflow to eliminate redundant
work and storage \citep{graphiler}; FeatGraph composes sparse traversal and
feature schedules \citep{featgraph}; Seastar generates fused forward/backward
kernels \citep{seastar}; and GALA composes graph-training transformations
\citep{gala}. Tiga likewise exploits message semantics, but additionally retains
generated relation constructors and explicit storage residence in its execution
model. These design differences are not comparative performance claims against
those compilers. GraphIt separates algorithms from schedules \citep{graphit},
while TACO and SparseTIR separate sparse iteration and storage choices
\citep{taco,sparsetir}. GraphChi and X-Stream establish that partitioned streaming
can make large graphs executable on a single machine \citep{graphchi,xstream}.
Tiga builds on these principles for differentiable numerical field programs;
neither fusion nor offloading is claimed as a new idea in isolation.
